% Page 43 — Alban : XVIème siècle (début)

\begin{center}
  {\Large\itshape Alban : XVI\textsuperscript{ème} siècle (début)}
\end{center}

\vspace{1em}

\lettrine{S}{i} les « de Moret » font parler d'eux en ce début de XVI\textsuperscript{ème} siècle, par le rang social qu'ils occupent, ainsi Etienne de Moret en 1545 est lieutenant du roi pour la baronnie d'Alban, d'autres familles vont aussi prendre importance. Un acte du 12 octobre 1537 fait état de « noble François Caudière », fils d'Antoine, d'Alban, vendant à Bernard Olm habitant la Mouline Ferran, un vignoble qu'il a à Malagousse (dans les parages de l'actuel Parrot). Mais c'est surtout la famille des « de Passemar » qui va s'assurer la main-mise sur le pays. Le 3 juillet 1561, Guillaume de Passemar, marchand d'Alban, achète à M. de Candols toutes les rentes de la paroisse de Massals et celles des masages (hameaux) d'Oulas, des Escombrazes, de la Rode, de la Bartabié, de Pomardelle, de Tels, de la Fanguezié, de la borie des Coubès, de la Tapie, de Margou, de la Besse, de Tessonière, des Estrébols, de Bordes, de Cantarrane, du mas del Bosc, des Thérondeis, le tout pour 600 livres. Dans le cadastre de 1598, nous verrons ce « noble Guillaume de Passemar », seigneur de Cantelauze, venir en tête de liste après Notre Dame de Piétat (N.D. d'Alban sous le vocable de N.D. de Pitié) après les Hospitaliers (successeurs des Templiers) après la Chapellenie de St Laurent fondée par les de Moret, mais avant Jacques de Moret.

\medskip

\lettrine{A}{u} civil comme au religieux, une réforme administrative est en train de s'opérer au commencement de ce XVI\textsuperscript{ème} siècle. Le comté de Castres a été directement rattaché à la couronne de France après la mort de Boffille de Juge. On peut donc dire d'une certaine façon qu'à ce moment de l'histoire, le roi François 1\textsuperscript{er} est seigneur baron d'Alban. En même temps, notre cité a pris rang de « chef lieu de Claverie », nous disons aujourd'hui tête de doyenné. La ville d'Alban a ainsi remplacé l'ancien archiprêtré de Négremont, dont le dernier archiprêtre a été démissionné en 1438. Cette « Claverie », notez qu'une ferme porte encore ce nom, vestige du passé, comprend 13 chapellenies (nous disons aujourd'hui : « fondations pieuses ») dont celle de la Cappélanyé dite de St Laurent dont les de Moret
